\subsection{DERIVATIVE OF THE INVERSE OF A FUNCTION}

PROPOSITION 4. Let E be a complete normed algebra with a unit element over R and let f be a function defined on an interval I ⊂ R, taking values in E, and differentiable at the point \(x_{0} \in I\) . If \(\mathbf{y}_{0} = \mathbf{f}(x_{0})\) is invertible \(^{2}\) in E, then the function \(x \mapsto (\mathbf{f}(x))^{-1}\) is defined on a neighbourhood of \(x_{0}\) (relative to I), and has a derivative equal to \(-\left(\mathbf{f}(x_{0})\right)^{-1}\mathbf{f}'(x_{0})\left(\mathbf{f}(x_{0})\right)^{-1}\) at \(x_{0}\) .

Indeed, the set of invertible elements in E is an open set on which the function \(y \mapsto y^{-1}\) is continuous (Gen.~Top., IX, p.~178); since f is continuous (relative to I) at \(x_{0}\) , \(\left(\mathbf{f}(x)\right)^{-1}\) is defined on a neighbourhood of \(x_{0}\) , and we have

\[
\left(\mathbf {f} (x)\right) ^ {- 1} - \left(\mathbf {f} \left(x _ {0}\right)\right) ^ {- 1} = \left(\mathbf {f} (x)\right) ^ {- 1} \left(\mathbf {f} \left(x _ {0}\right) - \mathbf {f} (x)\right) \left(\mathbf {f} \left(x _ {0}\right)\right) ^ {- 1}.
\]

The proposition thus follows from the continuity of \(y^{-1}\) on a neighbourhood of \(y_{0}\) and the continuity of xy on \(E \times E\) .

Examples. 1) The most important particular case is that where E is one of the fields R or C: if f is a function with real or complex values, differentiable at the point \(x_{0}\) , and such that \(f(x_{0}) \neq 0\) , then 1/f has derivative equal to \(-f'(x_{0})/(f(x_{0}))^{2}\) at \(x_{0}\) .

\begin{enumerate}
\def\labelenumi{\arabic{enumi})}
\setcounter{enumi}{1}
\tightlist
\item
  If \(U = (\alpha_{ij}(x))\) is a square matrix of order n, differentiable at \(x_{0}\) and invertible at this point, then \(U^{-1}\) has derivative equal to \(-U^{-1}U'U^{-1}\) at \(x_{0}\) .
\end{enumerate}

